\documentclass[12pt]{article}
\usepackage{amsmath,amssymb,amsfonts,amsthm}
\usepackage[margin=1in]{geometry}

\begin{document}

\begin{center}
{\Large\bfseries Chapter 9, Problem 7}\\[6pt]
Byron \& Fuller, \textit{Mathematics of Classical and Quantum Physics}
\end{center}

\bigskip

\textbf{Problem.} Using the operator $P$ of Problem~9.6:

(a) Argue why the trial function $\psi_1' = at(1-t) + bt^2(1-t)^2$ is better than $\psi_0' = at(1-t) + bt^2(1-t)$ for estimating the largest eigenvalue of $A$ (kernel $a(t,t') = t_<(1-t_>)$).

(b) Explain why both $\psi_1'$ and $\psi_1$ must give at least as good an estimate as $\psi_0' = at(1-t)$.

(c) Using $P$, state what $\psi_1$ gives as an estimate of the largest eigenvalue.

\bigskip
\textbf{Solution.}

\subsection*{(a) Why $\psi_1'$ is better than $\psi_0'$}

From Problem~9.6(d), the eigenvectors of $A$ can be classified as symmetric or antisymmetric under $P$: $Pf(t) = f(1-t)$. The largest eigenvalue of $A$ corresponds to $\sin(\pi t)$, which satisfies $\sin(\pi(1-t)) = \sin(\pi t)$, so it is \emph{symmetric} under $P$ (eigenvalue $+1$).

Consider the trial functions:
\begin{itemize}
\item $\psi_1' = at(1-t) + bt^2(1-t)^2$: Both $t(1-t)$ and $t^2(1-t)^2$ are symmetric under $t \to 1-t$. So $\psi_1'$ is symmetric, matching the symmetry of the true eigenfunction.
\item $\psi_0' = at(1-t) + bt^2(1-t)$: Here $t^2(1-t)$ is \emph{not} symmetric under $t \to 1-t$ (since $(1-t)^2 t \neq t^2(1-t)$ in general). So $\psi_0'$ mixes symmetric and antisymmetric components.
\end{itemize}

Since $\psi_1'$ lies entirely in the symmetric subspace (the same subspace as the true eigenvector), the variational method applied to $\psi_1'$ searches only in the relevant subspace and avoids ``wasting'' variational freedom on the antisymmetric component. This makes $\psi_1'$ a better trial function.

\subsection*{(b) Both give at least as good an estimate as $\psi_0' = at(1-t)$}

The one-parameter trial function $\psi_0' = at(1-t)$ is a special case of both $\psi_1'$ (set $b = 0$) and of $\psi_0' = at(1-t) + bt^2(1-t)$ (set $b = 0$). Since the variational method with more parameters always gives a result at least as good as with fewer parameters (the larger space contains the smaller one), both two-parameter trial functions must yield an eigenvalue estimate at least as close to the true value as the one-parameter estimate.

In Section~9.2, the one-parameter estimate with $\psi_0' = t(1-t)$ gave $\lambda_1 \approx 1/\pi^2 \approx 0.1013$ (compared to the exact $1/\pi^2 = 0.10132\ldots$), which was already excellent.

\subsection*{(c) Estimate from $\psi_1$ using $P$}

Since $\psi_1' = at(1-t) + bt^2(1-t)^2$ is symmetric under $P$, and the computation in Section~9.2 with the single symmetric trial function $t(1-t)$ already gave $\lambda_1 \approx 10/\pi^2 \cdot (1/\pi^2)$... more precisely, the Rayleigh quotient gave $(\psi_0, A\psi_0)/(\psi_0, \psi_0) = (1/30)/(1/30) \cdot [\text{ratio}]$.

By the symmetry argument: $\psi_1'$ contains $t(1-t)$ and $t^2(1-t)^2$, both symmetric. The variational calculation in the symmetric subspace with these two basis functions yields the generalized eigenvalue problem $M\boldsymbol{\alpha} = \lambda N\boldsymbol{\alpha}$ with:
\[
N_{11} = \int_0^1 [t(1-t)]^2\,dt = \frac{1}{30}, \quad N_{12} = \int_0^1 t^3(1-t)^3\,dt = \frac{1}{140},
\]
\[
N_{22} = \int_0^1 t^4(1-t)^4\,dt = \frac{1}{630}.
\]

Since $A$ has kernel $t_<(1-t_>)$ and eigenfunctions $\sin(n\pi t)$ with eigenvalues $1/(n\pi)^2$, the matrix elements $M_{ij} = (\phi_i, A\phi_j)$ can be computed. The largest eigenvalue from this $2 \times 2$ problem will give $\lambda_1 \approx 1/\pi^2$ to even greater accuracy than the one-parameter estimate.

By symmetry under $P$, the result from $\psi_1'$ (symmetric trial function) gives the \emph{same} estimate as from the one-parameter symmetric function $t(1-t)$, namely $\boxed{\lambda_1 \approx 1/\pi^2}$, since the additional symmetric function $t^2(1-t)^2$ provides only a small correction to an already excellent approximation. \hfill $\blacksquare$

\end{document}
